\section{Source substitutions and corrected density terms}

We follow \cite[proof of Theorem~5.2]{OpenAI2026PolynomialPEPS},
\texttt{04-compression.tex}, lines 279--383.
Let $W:\mathcal H_a\to\mathcal H_b$ be a chronological composition.
Its original source positions form the finite set $\mathcal E$.
For $e=(g,j)$, write $\Xi_g$ for the local labels of its owning gate,
$c_{g,\alpha}$ for the gate coefficients, and
$\nu_{e,\alpha}\in\C^{u_e}\otimes\C^{v_e}$ for its original source vector.
A replacement at $e$ may depend on $\alpha$, but not on labels of other gates.
All coordinate matrices below use fixed finite orthonormal bases on the
original input and output memories. The input matrix $\rho$ is arbitrary.
The finite-dimensional convention is that of the manuscript, lines 5--19;
no bound on private dimensions is imposed.

\begin{definition}[Substitution in the original composition]
    \label{def:peps_source_substitution}
    \lean{TNLean.PEPS.PairEffect.PreparedSourceGate.branchWithSources}
    \lean{TNLean.PEPS.PairEffect.SourceCircuit.evalWithSources}
    \lean{TNLean.PEPS.PairEffect.SourceCircuit.partialWithSources}
    \lean{TNLean.PEPS.PairEffect.SourceCircuit.selectedSourceVectors}
    \leanok
    \uses{def:peps_source_circuit,def:peps_prepared_source_gate,def:peps_partial_source_choices}
    Given vectors $\eta_{e,\alpha}$ in the original source halfspaces,
    replace the preparations in branch $\alpha$ of gate $g$ by these
    vectors, retain its remaining operations, and sum with its original
    coefficient $c_{g,\alpha}$. Compose these substituted gate operators
    in the original order; denote the result by $W(\eta)$.

    For affected parties $\mathcal A$, denote by
    $W_{\mathcal A,\xi}(\eta)$ the partial branch obtained by the same
    substitution at touched gates, retaining each untouched gate as its
    original aggregate operator. For $S\subseteq\mathcal E$, let
    $\nu[S\leftarrow\eta]$ equal $\eta$ on $S$ and $\nu$ elsewhere.
\end{definition}

\begin{theorem}[Summing untouched gates]
    \label{thm:peps_source_substitution_partial}
    \lean{TNLean.PEPS.PairEffect.SourceCircuit.evalWithSources_original}
    \lean{TNLean.PEPS.PairEffect.SourceCircuit.partialWithSources_original}
    \lean{TNLean.PEPS.PairEffect.SourceCircuit.evalWithSources_eq_sum_partialWithSources}
    \lean{TNLean.PEPS.PairEffect.SourceCircuit.eval_selectedSourceVectors_eq_sum}
    \leanok
    \uses{def:peps_source_substitution}
    Original vectors recover $W(\nu)=W$ and the original partial branches.
    Suppose $\eta_{e,\alpha}=\nu_{e,\alpha}$ whenever the owning gate
    of $e$ does not touch $\mathcal A$. If $J_a,J_b$ are the canonical
    owner-identification isometries, then
    \begin{align}
        J_b W(\eta)
         & =\sum_{\xi}c_{\mathcal A,\xi}
        W_{\mathcal A,\xi}(\eta)J_a.
        \label{eq:peps_substitution_partial}
    \end{align}
    In particular, (\ref{eq:peps_substitution_partial}) holds for
    $\nu[S\leftarrow\eta]$ if every position in $S$ has an endpoint
    in $\mathcal A$. The substituted vectors need not be normalized.
\end{theorem}

\begin{proof}
    \lean{TNLean.PEPS.PairEffect.Layout.affectedOwner_mapOwner_eq_none}
    \leanok
    \uses{thm:peps_partial_source_evaluation,thm:peps_source_circuit_incidence}
    Induct on the chronological composition. At a touched gate, the
    sum is its defining branch sum. At an untouched gate all source
    vectors are original, and every participant has exterior owner;
    its branch sum is therefore its original aggregate operator.
    Composition distributes through the finite sums, while idle
    registers tensor each term with an identity. Every endpoint of a
    selected source belongs to its owning gate, which proves the
    final specialization.
\end{proof}

\begin{definition}[The canonical remaining operations]
    \label{def:peps_canonical_source_residual}
    \lean{TNLean.PEPS.PairEffect.Word.residualInSlots}
    \lean{TNLean.PEPS.PairEffect.SourceCircuit.partialResidual}
    \leanok
    \uses{thm:peps_source_preparation,thm:peps_partial_fixed_layout}
    Moving the actual sources of a word to its initial registers
    leaves a specified composition of its original remaining operations.
    Identify its source-register list with any equal prescribed list.
    Denote the resulting composition by $V$.
    For the partial branch $W_{\mathcal A,\xi}$, use the fixed retained
    positions $\mathcal E_{\mathcal A}$ and write $V_{\mathcal A,\xi}$.
\end{definition}

\begin{theorem}[Recovery by the canonical remaining operations]
    \label{thm:peps_canonical_source_recovery}
    \lean{TNLean.PEPS.PairEffect.Word.sources_residualInSlots}
    \lean{TNLean.PEPS.PairEffect.Word.isAllowed_residualInSlots}
    \lean{TNLean.PEPS.PairEffect.Word.eval_residualInSlots_prepare}
    \lean{TNLean.PEPS.PairEffect.SourceCircuit.slotLayout_partialSlots_eq}
    \lean{TNLean.PEPS.PairEffect.SourceCircuit.sources_partialResidual}
    \lean{TNLean.PEPS.PairEffect.SourceCircuit.isAllowed_partialResidual}
    \lean{TNLean.PEPS.PairEffect.SourceCircuit.eval_partialResidual_prepare}
    \leanok
    \uses{def:peps_canonical_source_residual}
    The composition $V$ has no pair sources. It is allowed if the
    original word is allowed. Whenever a preparation in the prescribed
    positions has exactly the original ordered source inventory,
    $VP_\eta$ is the original word operator.
    Consequently the allowed partial branches satisfy
    \begin{align}
        V_{\mathcal A,\xi}P_{\nu_\xi}
         & =W_{\mathcal A,\xi},
         & (\nu_\xi)_i          & =\nu_{e_i,\xi_{g_i}}.
        \label{eq:peps_canonical_recovery}
    \end{align}
\end{theorem}

\begin{proof}
    \leanok
    \uses{thm:peps_source_preparation,thm:peps_partial_original_source_vector}
    Transport the source-preparation identity along the equality of
    ordered register lists. The inventory at index $i$ is the original
    vector selected by $\xi_{g_i}$, which gives
    (\ref{eq:peps_canonical_recovery}).
\end{proof}

\begin{theorem}[Uniform recovery after source substitution]
    \label{thm:peps_canonical_substitution_recovery}
    \lean{TNLean.PEPS.PairEffect.SourceCircuit.layout_sources_partialWithSources}
    \lean{TNLean.PEPS.PairEffect.SourceCircuit.residualInSlots_partialWithSources}
    \lean{TNLean.PEPS.PairEffect.SourceCircuit.ofSlots_partialWithSources}
    \lean{TNLean.PEPS.PairEffect.SourceCircuit.eval_partialResidual_prepare_substituted}
    \leanok
    \uses{def:peps_source_substitution,def:peps_canonical_source_residual}
    For arbitrary $\eta$, the ordered pair-source inventory of
    $W_{\mathcal A,\xi}(\eta)$ has the same retained positions and
    has vector $\eta_{e_i,\xi_{g_i}}$ at position $i$.
    If every source with both endpoints outside $\mathcal A$ retains
    its original vector, the canonical remaining composition is
    literally $V_{\mathcal A,\xi}$, independently of the retained
    vectors. Thus
    \begin{align}
        V_{\mathcal A,\xi}P_{\eta_\xi}
         & =W_{\mathcal A,\xi}(\eta),
         & (\eta_\xi)_i               & =\eta_{e_i,\xi_{g_i}}.
        \label{eq:peps_substituted_recovery}
    \end{align}
\end{theorem}

\begin{proof}
    \leanok
    \uses{thm:peps_partial_source_order,thm:peps_canonical_source_recovery}
    Induction preserves the complete ordered source records.
    Surviving pair preparations change only their vectors. A source
    whose endpoints are both exterior becomes a local preparation;
    the stated condition keeps that local operation unchanged.
    All other remaining operations agree. Apply the canonical recovery
    identity to the resulting exact inventory.
\end{proof}

\begin{theorem}[Weighted matrix sums]
    \label{thm:peps_weighted_density_sum}
    \lean{Matrix.sum_smul_mul_mul_conjTranspose}
    \leanok
    For finite coefficient families $c_\alpha,d_\beta$, rectangular
    matrices $K_\alpha,L_\beta$, and a compatible rectangular matrix $R$,
    \begin{align}
        \Bigl(\sum_\alpha c_\alpha K_\alpha\Bigr)R
        \Bigl(\sum_\beta d_\beta L_\beta\Bigr)^\dagger
         & =\sum_{\alpha,\beta}c_\alpha\overline{d_\beta}
        K_\alpha R L_\beta^\dagger.
        \label{eq:peps_weighted_density_sum}
    \end{align}
\end{theorem}
\begin{proof}
    \leanok
    Distribute the products over the finite sums and conjugate the
    coefficients in the adjoint.
\end{proof}

\begin{definition}[Corrected terms at original source positions]
    \label{def:peps_original_corrected_term}
    \lean{TNLean.PEPS.PairEffect.SourceCircuit.SourceBasisChoice}
    \lean{TNLean.PEPS.PairEffect.SourceCircuit.sourceBasisChoiceFintype}
    \lean{TNLean.PEPS.PairEffect.SourceCircuit.labelBasisVectors}
    \lean{TNLean.PEPS.PairEffect.SourceCircuit.sourceBasisMatrix}
    \lean{TNLean.PEPS.PairEffect.SourceCircuit.correctedSourceTerm}
    \leanok
    \uses{def:peps_source_substitution}
    Fix $S\subseteq\mathcal E$. A choice $u$ specifies at each $e=(g,j)\in S$
    a local label $\alpha_e\in\Xi_g$ and an endpoint basis index $a_e$.
    In $W$, replace the vector at $(e,\gamma)$ by the basis tensor
    indexed by $a_e$ if $\gamma=\alpha_e$, and by zero otherwise;
    retain original vectors outside $S$. Let $M_{S,u}$ be the resulting
    circuit matrix. For arbitrary source matrices $E_{e,\alpha,\beta}$,
    define
    \begin{align}
        F_S(E;\rho)
         & :=\sum_{u=(\alpha,a),\,v=(\beta,b)}
        \left(\prod_{e\in S}E_{e,\alpha_e,\beta_e}(a_e,b_e)\right)
        M_{S,u}\rho M_{S,v}^{\dagger}.
        \label{eq:peps_original_corrected_term}
    \end{align}
    This definition uses the original positions and local labels,
    before any affected set is chosen.
\end{definition}

\begin{theorem}[The empty corrected set]
    \label{thm:peps_corrected_empty}
    \lean{TNLean.PEPS.PairEffect.SourceCircuit.correctedSourceTerm_empty}
    \leanok
    \uses{def:peps_original_corrected_term}
    If $M$ is the original circuit matrix, then $F_\varnothing(E;\rho)=M\rho M^\dagger$.
\end{theorem}
\begin{proof}
    \leanok
    \uses{thm:peps_source_substitution_partial}
    There is one assignment on the empty set, its product is one,
    and every source remains original.
\end{proof}

\begin{definition}[Partially expanded basis substitutions]
    \label{def:peps_partial_basis_substitution}
    \lean{TNLean.PEPS.PairEffect.SourceCircuit.partialSourceBasisMatrix}
    \lean{TNLean.PEPS.PairEffect.SourceCircuit.SelectedSourceCoordinates}
    \lean{TNLean.PEPS.PairEffect.SourceCircuit.selectedSourceCoordinatesFintype}
    \lean{TNLean.PEPS.PairEffect.SourceCircuit.coordinateBasisVectors}
    \lean{TNLean.PEPS.PairEffect.SourceCircuit.partialCoordinateBasisMatrix}
    \leanok
    \uses{def:peps_original_corrected_term,def:peps_partial_source_choices}
    Let $M_{\mathcal A,\xi;S,u}$ be the matrix of the corresponding
    substituted partial branch, transported back to the original
    input and output coordinates by $J_a,J_b$.
    A coordinate choice $a$ on $S$ specifies only the endpoint basis
    indices. Write $K_{\mathcal A,\xi;S,a}$ for the same partial matrix
    with these basis tensors at selected positions, without imposing
    additional local-label conditions.
\end{definition}

\begin{theorem}[Partial expansion of the same corrected term]
    \label{thm:peps_corrected_partial}
    \lean{TNLean.PEPS.PairEffect.SourceCircuit.sourceBasisMatrix_eq_sum_partial}
    \lean{TNLean.PEPS.PairEffect.SourceCircuit.correctedSourceTerm_eq_sum_partial}
    \lean{TNLean.PEPS.PairEffect.SourceCircuit.partialSourceBasisMatrix_eq_ite}
    \lean{TNLean.PEPS.PairEffect.SourceCircuit.correctedSourceTerm_eq_sum_coordinates}
    \lean{TNLean.PEPS.PairEffect.SourceCircuit.partialCoordinateBasisMatrix_eq_preparedMatrix}
    \leanok
    \uses{def:peps_partial_basis_substitution}
    Suppose every $e\in S$ has an endpoint in $\mathcal A$.
    Then $M_{S,u}=\sum_\xi c_{\mathcal A,\xi}M_{\mathcal A,\xi;S,u}$.
    A summand $M_{\mathcal A,\xi;S,u}$ equals
    $K_{\mathcal A,\xi;S,a}$ if $\alpha_e=\xi_{g_e}$ at every
    selected position, and is zero otherwise. Consequently
    \begin{align}
        F_S(E;\rho)
         & =\sum_{\xi,\zeta}c_{\mathcal A,\xi}\overline{c_{\mathcal A,\zeta}}
        \sum_{a,b}
        \left(\prod_{e\in S}
        E_{e,\xi_{g_e},\zeta_{g_e}}(a_e,b_e)\right)
        K_{\mathcal A,\xi;S,a}\rho K_{\mathcal A,\zeta;S,b}^{\dagger}.
        \label{eq:peps_corrected_partial}
    \end{align}
    Each $K$ is the matrix of the canonical remaining composition
    applied to the indicated basis preparations, with the original
    locally selected vectors at every unselected retained position.
\end{theorem}

\begin{proof}
    \leanok
    \uses{thm:peps_source_substitution_partial,thm:peps_weighted_density_sum,
        thm:peps_canonical_substitution_recovery,thm:peps_source_preparation_sums}
    Apply (\ref{eq:peps_substitution_partial}) to each basis substitution
    and distribute the ket and bra sums using
    (\ref{eq:peps_weighted_density_sum}). If a selected label disagrees
    with the branch label, its source vector is zero. Multilinearity
    of preparation makes the corresponding matrix zero. Otherwise the
    substitution is exactly the prescribed basis preparation, by
    (\ref{eq:peps_substituted_recovery}). Summing the local-label
    conditions leaves (\ref{eq:peps_corrected_partial}). Selected
    positions at the same gate impose the same label; the coefficient
    of that gate still occurs only once in $c_{\mathcal A,\xi}$.
\end{proof}

\begin{definition}[Selected retained positions]
    \label{def:peps_selected_retained_positions}
    \lean{TNLean.PEPS.PairEffect.SourceCircuit.selectedPartialSlots}
    \lean{TNLean.PEPS.PairEffect.SourceCircuit.mem_selectedPartialSlots}
    \lean{TNLean.PEPS.PairEffect.SourceCircuit.selectedPartialSlotEquiv}
    \lean{TNLean.PEPS.PairEffect.SourceCircuit.selectedSourceCoordinateEquiv}
    \lean{TNLean.PEPS.PairEffect.SourceCircuit.selectedSourceCoordinateEquiv_apply}
    \leanok
    \uses{def:peps_partial_fixed_positions}
    Let $i\mapsto e_i$ enumerate the retained positions. Set
    $S_{\mathcal A}=\{i:e_i\in S\}$. If every selected position has an
    affected endpoint, $i\mapsto e_i$ restricts to a bijection
    $S_{\mathcal A}\to S$. Precomposition with this bijection identifies
    the two families of endpoint-coordinate assignments, point by point.
\end{definition}
\begin{theorem}[Changing the enumeration of selected sources]
    \label{thm:peps_selected_retained_coordinates}
    \lean{TNLean.PEPS.PairEffect.SourceCircuit.partialCoordinateBasisMatrix_reindex}
    \leanok
    \uses{def:peps_selected_retained_positions,thm:peps_canonical_substitution_recovery}
    Under this bijection, $K_{\mathcal A,\xi;S,a}$ is exactly the
    prepared matrix of $V_{\mathcal A,\xi}$ with the corresponding
    basis vectors at $S_{\mathcal A}$ and original source vectors elsewhere.
\end{theorem}
\begin{proof}
    \leanok
    Compare the vector at each retained index $i$. Membership in
    $S_{\mathcal A}$ is precisely $e_i\in S$, and both expressions
    select the same coordinate in that case.
\end{proof}

For a fixed ordered source list, let $K_V(a)$ be the matrix obtained
from a remaining composition $V$ by preparing its endpoint vectors
indexed by $a$. For ket and bra endpoint families, put
$H_{V,V'}(a,b)=K_V(a)\rho K_{V'}(b)^\dagger$ and write
\begin{align}
    \mathcal C_H(Y)
     & :=\sum_{a,b}\left(\prod_iY_i(a_i,b_i)\right)H(a,b).
    \label{eq:peps_source_contraction}
\end{align}
The endpoint families on the two sides may have different finite
index sets.

\begin{theorem}[Prepared densities as source contractions]
    \label{thm:peps_prepared_source_contraction}
    \lean{TNLean.PEPS.PairEffect.Word.sourceContraction_preparedDensityCoefficient}
    \lean{TNLean.PEPS.PairEffect.Word.sourceContraction_preparedDensityCoefficient_basis}
    \leanok
    \uses{thm:peps_prepared_density_components}
    Suppose $\eta_i=\sum_a r_i(a)u_i(a_1)\otimes v_i(a_2)$ and
    $\zeta_i=\sum_b s_i(b)u'_i(b_1)\otimes v'_i(b_2)$.
    With $X_i=r_i s_i^\dagger$, the actual prepared matrices satisfy
    \begin{align}
        \mathcal C_{H_{V,V'}}(X)
         & =K_V(\eta)\rho K_{V'}(\zeta)^\dagger.
        \label{eq:peps_prepared_contraction}
    \end{align}
    No spanning or normalization condition is needed for these finite
    expansions. In particular the identity holds for the coordinate
    vectors of arbitrary sources in independently chosen endpoint
    orthonormal bases.
\end{theorem}
\begin{proof}
    \leanok
    Pair the ket and bra assignments at each position. Expanding their
    rank-one matrix entries gives the product of the ket coefficients
    and conjugated bra coefficients. The finite preparation expansion
    then gives (\ref{eq:peps_prepared_contraction}).
\end{proof}

\begin{definition}[Replacement of source operators]
    \label{def:peps_prepared_source_replacement}
    \lean{TNLean.PEPS.PairEffect.Word.sourceOperatorDensity}
    \leanok
    For actual remaining compositions $V_\xi$ and weights $c_\xi$,
    define
    \begin{align}
        \mathcal R(Y)
         & :=\sum_{\xi,\zeta}c_\xi\overline{c_\zeta}
        \mathcal C_{H_{V_\xi,V_\zeta}}(Y_{\xi,\zeta}).
        \label{eq:peps_prepared_replacement}
    \end{align}
    Source matrices may depend on both branch labels; ket and bra
    endpoint-coordinate sets may vary with their respective branch labels.
    The matrices need not be positive or Hermitian.
\end{definition}
\begin{theorem}[The prepared corrected-subset identity]
    \label{thm:peps_prepared_subset_identity}
    \lean{TNLean.PEPS.PairEffect.Word.sourceOperatorDensity_sub}
    \lean{TNLean.PEPS.PairEffect.Word.sourceOperatorDensity_exact}
    \lean{TNLean.PEPS.PairEffect.Word.sourceOperatorDensity_sub_prepared}
    \leanok
    \uses{def:peps_prepared_source_replacement}
    If $X$ is any family of source matrices, then
    \begin{align}
        \mathcal R(X+E)-\mathcal R(X)
         & =\sum_{\varnothing\ne S\subseteq\{0,\ldots,r-1\}}
        \mathcal R(E\text{ on }S,\ X\text{ off }S).
        \label{eq:peps_prepared_subsets}
    \end{align}
    For exact rank-one coordinate matrices of sources $\eta_\xi$,
    $\mathcal R(X)=M\rho M^\dagger$, where
    $M=\sum_\xi c_\xi K_{V_\xi}(\eta_\xi)$.
\end{theorem}
\begin{proof}
    \leanok
    \uses{thm:peps_prepared_source_contraction,thm:peps_weighted_density_sum}
    Expand the product in (\ref{eq:peps_source_contraction}) over subsets
    and subtract its empty-subset term. Apply
    (\ref{eq:peps_prepared_contraction}) and
    (\ref{eq:peps_weighted_density_sum}) for the exact-source assertion.
\end{proof}

\begin{theorem}[Expanding only the selected source matrices]
    \label{thm:peps_selected_source_contraction}
    \lean{TNLean.PEPS.PairEffect.Word.sourceContraction_preparedDensityCoefficient_selected}
    \leanok
    \uses{thm:peps_prepared_source_contraction}
    Let $X_i=\ketbra{\eta_i}{\zeta_i}$ in the chosen ket and bra
    endpoint bases. For any set $S$ of positions,
    \begin{align}
        \mathcal C_{H_{V,V'}}(E\text{ on }S,\ X\text{ off }S)
         & =\sum_{a,b}\left(\prod_{i\in S}E_i(a_i,b_i)\right)
        K_V(\eta[S\leftarrow a])\rho
        K_{V'}(\zeta[S\leftarrow b])^\dagger.
        \label{eq:peps_selected_contraction}
    \end{align}
    Here $a,b$ range only over endpoint-coordinate assignments on $S$,
    and the notation on the right replaces each selected vector by
    its endpoint basis tensor. Ket and bra bases may differ.
\end{theorem}
\begin{proof}
    \leanok
    Expand each selected matrix in matrix units. Multilinearity of
    (\ref{eq:peps_source_contraction}) leaves a sum over only those
    positions. A selected matrix unit is the rank-one coordinate
    matrix of its ket and bra basis tensors. Apply
    (\ref{eq:peps_prepared_contraction}) to each term.
\end{proof}

\begin{theorem}[The corrected term in a partial expansion]
    \label{thm:peps_corrected_source_contraction}
    \lean{TNLean.PEPS.PairEffect.SourceCircuit.correctedSourceTerm_eq_sourceContraction}
    \leanok
    \uses{def:peps_original_corrected_term,def:peps_selected_retained_positions}
    Suppose every $e\in S$ has an endpoint in $\mathcal A$.
    In the common retained positions, set
    \begin{align}
        Y_{\xi,\zeta,i}
         & :=\begin{cases}
                 E_{e_i,\xi_{g_i},\zeta_{g_i}},                       & e_i\in S,    \\
                 \ketbra{\nu_{e_i,\xi_{g_i}}}{\nu_{e_i,\zeta_{g_i}}}, & e_i\notin S.
             \end{cases}
        \label{eq:peps_partial_corrected_sources}
    \end{align}
    The same globally defined term then satisfies
    \begin{align}
        F_S(E;\rho)
         & =\sum_{\xi,\zeta}c_{\mathcal A,\xi}\overline{c_{\mathcal A,\zeta}}
        \mathcal C_{H_{V_{\mathcal A,\xi},V_{\mathcal A,\zeta}}}
        (Y_{\xi,\zeta}).
        \label{eq:peps_corrected_contraction}
    \end{align}
\end{theorem}
\begin{proof}
    \leanok
    \uses{thm:peps_corrected_partial,thm:peps_selected_retained_coordinates,
        thm:peps_selected_source_contraction}
    Apply (\ref{eq:peps_selected_contraction}) to each partial ket--bra
    pair and reindex its selected coordinates along the bijection
    $S_{\mathcal A}\to S$. The resulting expression is exactly
    (\ref{eq:peps_corrected_partial}).
\end{proof}

\begin{definition}[All original source positions]
    \label{def:peps_all_original_source_slots}
    \lean{TNLean.PEPS.PairEffect.SourceCircuit.allSourceSlotEquiv}
    \lean{TNLean.PEPS.PairEffect.SourceCircuit.allSourceSlotEquiv_apply}
    \leanok
    \uses{def:peps_partial_fixed_positions}
    Taking every party to be affected retains every original source
    position. Its ordered enumeration is a bijection
    $\{0,\ldots,r-1\}\to\mathcal E$. No finiteness assumption on the
    set of parties is needed.
\end{definition}
\begin{theorem}[Reindexing nonempty source subsets]
    \label{thm:peps_all_source_subsets}
    \lean{TNLean.PEPS.PairEffect.SourceCircuit.sum_nonempty_allSourceSlots}
    \leanok
    \uses{def:peps_all_original_source_slots}
    This bijection identifies sums over nonempty subsets of the
    retained indices with sums over nonempty subsets of $\mathcal E$,
    for every function valued in a commutative additive monoid.
\end{theorem}
\begin{proof}
    \leanok
    The induced bijection on finite subsets preserves emptiness.
\end{proof}

\begin{definition}[Replacement in the complete circuit density]
    \label{def:peps_global_source_replacement}
    \lean{TNLean.PEPS.PairEffect.SourceCircuit.exactSourceMatrix}
    \lean{TNLean.PEPS.PairEffect.SourceCircuit.sourceReplacedDensity}
    \leanok
    \uses{def:peps_prepared_source_replacement,def:peps_all_original_source_slots}
    Put $X_{e,\alpha,\beta}=\ketbra{\nu_{e,\alpha}}{\nu_{e,\beta}}$.
    To define $\widehat\rho(Y)$, take all parties to be affected in
    (\ref{eq:peps_prepared_replacement}), use the actual canonical
    remaining compositions, and assign at position $e_i$ the matrix
    $Y_{e_i,\xi_{g_i},\zeta_{g_i}}$. Thus $\widehat\rho(Y)$ is defined
    directly by local source-operator replacement, independently of
    any corrected-subset formula.
\end{definition}
\begin{theorem}[The global corrected-subset identity]
    \label{thm:peps_global_corrected_subsets}
    \lean{TNLean.PEPS.PairEffect.SourceCircuit.sourceReplacedDensity_exact}
    \lean{TNLean.PEPS.PairEffect.SourceCircuit.sourceReplacedDensity_piecewise}
    \lean{TNLean.PEPS.PairEffect.SourceCircuit.sourceReplacedDensity_sub_exact}
    \leanok
    \uses{def:peps_global_source_replacement}
    The exact sources recover $\widehat\rho(X)=M\rho M^\dagger$.
    Replacing precisely the original positions in $S$ gives
    $\widehat\rho(E\text{ on }S,X\text{ off }S)=F_S(E;\rho)$.
    Consequently
    \begin{align}
        \widehat\rho(X+E)-M\rho M^\dagger
         & =\sum_{\varnothing\ne S\subseteq\mathcal E}F_S(E;\rho).
        \label{eq:peps_global_corrected_subsets}
    \end{align}
\end{theorem}
\begin{proof}
    \leanok
    \uses{thm:peps_corrected_source_contraction,thm:peps_corrected_empty,
        thm:peps_prepared_subset_identity,thm:peps_all_source_subsets}
    Apply (\ref{eq:peps_corrected_contraction}) with every party
    affected, first to the empty selected set and then to arbitrary
    $S$. This gives the two recovery identities. Expand the source
    operators by (\ref{eq:peps_prepared_subsets}) and reindex the
    nonempty subsets by their original occurrences.
\end{proof}

\begin{definition}[Physical output after discarding retained registers]
    \label{def:peps_physical_source_replacement}
    \lean{TNLean.PEPS.PairEffect.SourceCircuit.physicalDensity}
    \lean{TNLean.PEPS.PairEffect.SourceCircuit.physicalCorrectedSourceTerm}
    \lean{TNLean.PEPS.PairEffect.SourceCircuit.physicalSourceReplacedDensity}
    \leanok
    \uses{def:peps_global_source_replacement,def:peps_original_corrected_term}
    Suppose the final ordered register list is the physical list
    followed by the discarded list. Choose finite orthonormal bases
    of both memories and their tensor-product basis on the full
    output, transported through the canonical append isometry.
    Define
    \begin{align}
        \rho_{\mathrm{phys}}            & :=\operatorname{Tr}_{\mathrm{disc}}(M\rho M^\dagger), \\
        \widehat\rho_{\mathrm{phys}}(Y) & :=\operatorname{Tr}_{\mathrm{disc}}\widehat\rho(Y),   \\
        F_S^{\mathrm{phys}}(E;\rho)     & :=\operatorname{Tr}_{\mathrm{disc}}F_S(E;\rho).
    \end{align}
\end{definition}
\begin{theorem}[The physical corrected-subset identity]
    \label{thm:peps_physical_corrected_subsets}
    \lean{TNLean.PEPS.PairEffect.SourceCircuit.physicalSourceReplacedDensity_exact}
    \lean{TNLean.PEPS.PairEffect.SourceCircuit.physicalSourceReplacedDensity_sub_exact}
    \lean{TNLean.PEPS.PairEffect.SourceCircuit.rectangularTraceNorm_physicalSourceReplacedDensity_sub_le}
    \leanok
    \uses{def:peps_physical_source_replacement}
    Exact sources recover $\widehat\rho_{\mathrm{phys}}(X)=\rho_{\mathrm{phys}}$.
    The physical replacement error satisfies
    \begin{align}
        \widehat\rho_{\mathrm{phys}}(X+E)-\rho_{\mathrm{phys}}
         & =\sum_{\varnothing\ne S\subseteq\mathcal E}F_S^{\mathrm{phys}}(E;\rho),
        \label{eq:peps_physical_corrected_subsets}                                 \\
        \bigl\|\widehat\rho_{\mathrm{phys}}(X+E)-\rho_{\mathrm{phys}}\bigr\|_1
         & \leq\sum_{\varnothing\ne S\subseteq\mathcal E}
        \bigl\|F_S^{\mathrm{phys}}(E;\rho)\bigr\|_1.
    \end{align}
\end{theorem}
\begin{proof}
    \leanok
    \uses{thm:peps_global_corrected_subsets}
    Apply the linear partial trace to
    (\ref{eq:peps_global_corrected_subsets}), then use the triangle
    inequality for trace norm. No discarded-memory dimension factor
    appears in either assertion.
\end{proof}
